\chapter{Background and Literature Review}\label{chap:litreview}

\section{Nix Foundations}

This section introduces the Nix concepts required to understand why Nix is the choice of . Nix is both a package manager and a language for describing how software and development environments are obtained. Unlike an imperative installation procedure, which modifies a shared system state step by step, a Nix expression describes the inputs from which a package or environment should be produced. Nix then evaluates that description and creates a derivation. This distinction is central to the thesis: an agent generates a declaration rather than a sequence of host-specific installation commands.

\subsection{Declarative Packages and Nixpkgs}

In Nix, a \emph{derivation} is a build specification. It identifies the builder, its source inputs, build-time and run-time dependencies, configuration arguments, and expected outputs. Nix evaluates derivations as part of a dependency graph, allowing a package to depend on specific versions of compilers, libraries, and build tools rather than on whichever versions happen to be installed globally. The functional model aims to make a build result a function of its declared inputs, thereby avoiding hidden dependencies on mutable global directories~\cite{dolstraNixSafePolicyFree,dolstraPurelyFunctionalSystema}.

The collection of such package definitions is nixpkgs. It exposes packages as named attributes, such as \texttt{python3}, \texttt{cmake}, or \texttt{rustc}, and parameterizes many definitions by the target platform. A generated development environment therefore does not install packages by conventional system mutation; it selects attributes from a particular Nixpkgs revision and composes them into a shell or package output. This is why an invalid attribute name, a deprecated package, or a package unavailable on the target platform can prevent an otherwise plausible environment from evaluating. RQ1 evaluates whether the agent can make these selections correctly.

Declarative does not mean that every result is permanently identical in every context. Reproducibility depends on retaining the relevant input revisions, source artifacts, platform assumptions, and configuration. It also depends on the continued ability to obtain or build the required dependency closure. Consequently, the explicit declarations produced by Nix provide a strong basis for reproducibility, but they do not by themselves guarantee that a new student machine can activate an environment within an acceptable time.

\subsection{The Nix Store and Dependency Closures}

Nix places realized build outputs in the Nix store, conventionally \texttt{/nix/store}. Store paths are named with a cryptographic hash followed by a human-readable package name, for example \texttt{/nix/store/<hash>-gcc-<version>}. The hash distinguishes the declared input configuration associated with that output. Store paths are treated as immutable once created, so an upgrade does not overwrite a library used by an existing environment; instead, it creates or selects a distinct store path. This permits multiple versions of a tool or library to coexist without sharing a mutable installation prefix~\cite{dolstraNixOSPurelyFunctional2010}.

Every realized package refers to other store paths that it requires at run time. The transitive set of those paths is its \emph{closure}. A development environment is therefore more than a short list of executables: it includes the compiler, libraries, interpreters, shell hooks, and other transitive dependencies needed by those executables. Nix makes this closure explicit and can construct a shell in which selected package binaries are available through \texttt{PATH}. The command \texttt{nix develop} is the primary interface used in this thesis to activate such a shell.

The store model explains both the benefit and the limitation evaluated by RQ2. A host on which an environment has already been used may contain a large, warmed store whose paths mask missing declarations or make activation appear faster than it would be for a student. The clean-container replay therefore starts each validation stage without a mounted host store or project directory. It tests whether the generated declaration and its transitive closure can be obtained and used independently of the authoring machine's prior state.

\subsection{Flakes, Locked Inputs, and Development Shells}

A flake is a convention for organizing a Nix project around two files. The \texttt{flake.nix} file declares inputs and outputs. Inputs commonly include Nixpkgs, the nix package manager or other Git repositories; outputs may include packages, applications, NixOS configurations, or development shells. In this thesis, the relevant output is normally a \texttt{devShell}, which lists the tools and environment configuration required by a course assignment. A student activates that output with the \texttt{nix develop} shell command.

The companion \texttt{flake.lock} file records resolved input revisions and associated metadata. It is a commitment to a specific dependency graph: once an instructor reviews and commits it, students use the same pinned inputs rather than silently receiving the newest revision of Nixpkgs. Updating an input with \texttt{nix flake update} is therefore a maintenance action that intentionally changes the lock file; it is not part of ordinary activation of a fixed course environment. Pinning supports repeatability across student machines, but a lock file is not an offline archive. If an input or store path is absent locally, Nix may still need network access to retrieve it.

\begin{lstlisting}[caption={\texttt{flake.lock} \texttt{nixpkgs} Attributes}, label={lst:flake-lock-nixpkgs-attributes}]
  "nixpkgs": {
    "locked": {
      "lastModified": 1791011291,
      "narHash": "sha256-V6OBJccNKBo/ktknxlJXZChPoI0NE2lembArBO3CE9k=",
      "owner": "NixOS",
      "repo": "nixpkgs",
      "rev": "73e728ddb6b7a12d18808f510813a13ee1fe4cce",
      "type": "github"
    },
  }
\end{lstlisting}

\texttt{revision} is the most recent commit

This distinction is important to RQ3. An outage affecting a source host such as GitHub can interfere with updating a flake and can also interfere with first use when a locked input has not yet been fetched. These cases should not be conflated with failure of an already available, locked environment. Likewise, an instructor's decision to refresh a lock file can reveal renamed, removed, or insufficiently maintained package attributes that were not present in the version used for a prior course offering.

\subsection{Realization, Binary Caches, and Trust}

Nix evaluation determines whether an expression is syntactically and semantically valid and identifies the derivations it requires. Realization obtains the corresponding store paths. During realization, Nix first attempts to substitute a pre-built store path from a configured binary cache; the public cache at \texttt{cache.nixos.org} is the customary default. If no acceptable substitute is available, Nix can build the derivation from source, provided that its sources and build dependencies can be obtained. This fallback improves completeness but may require substantial time and local computing resources for large dependency closures.

A \emph{substituter} is an additional binary-cache endpoint. It can reduce setup time by serving store paths that are absent from the primary cache, including course-specific builds maintained by an institution. However, cache use is also a supply-chain trust decision. Nix accepts binary substitutes according to configured public signing keys; adding a third-party cache together with its key expands the set of entities authorized to provide binaries that students execute. Thus, a substitute may improve availability while increasing the trusted computing base~\cite{dolstraNixSafePolicyFree}.

These mechanisms establish the technical basis for RQ3. Classroom reliability depends on more than whether a \texttt{flake.nix} evaluates: locked inputs must remain obtainable when needed, package definitions must remain suitable for the course, required binaries must be available or buildable within the instructional time budget, and any cache used to improve availability must be explicitly trusted. The controlled reliability conditions in Chapter~\ref{chap:testing} measure these properties separately.

\section{Related Work}

\subsection{Standardizing Educational Environments via Virtualization and Containers}

Standardizing student development environments has remained a persistent pedagogical challenge in computer science education~\cite{stokerUsingVirtualMachines2013}. Historically, institutions attempted to resolve configuration discrepancies by deploying centralized compute clusters or multi-user Linux servers accessible via SSH or remote desktop protocols. Valstar et~al.~\cite{valstarUsingDevContainersStandardize2020} critically analyze the friction inherent in these pre-configured central servers, noting that instructors must continually coordinate with central IT personnel who manage infrastructure across disparate academic departments. This administrative bottleneck results in sluggish response times when course toolchains require updates during active semesters.

Furthermore, centralized remote architectures impose rigid network requirements~\cite{valstarUsingDevContainersStandardize2020}. Students must maintain persistent, low-latency network connections, either on campus or through virtual private networks (VPNs) provided by the institution. Unstable internet connections introduce substantial input latency and unexpected disconnects, severely disrupting the student learning experience during programming assignments and interactive debugging sessions. To alleviate these network and administrative dependencies, educators shifted toward client-side containerization through Development Containers (DevContainers). DevContainers leverage Docker to package compilers, runtimes, and editor extensions into a local container specification, ensuring that every enrolled student executes code in an identical software sandbox~\cite{valstarUsingDevContainersStandardize2020}. However, this approach shifts the operational burden to the student's personal hardware, demanding disk storage and memory capacity that student laptops may lack.

\subsection{Synthesis of Infrastructure as Code Using Large Language Models}

The rapid proliferation of Large Language Models has spurred widespread investigation into automating Infrastructure as Code (IaC) and system configuration management~\cite{kateGenerativeAIInfrastructure}. Rather than manually composing complex scripts, developers can specify deployment goals in natural language and delegate script generation to automated models. Srivatsa et~al.~\cite{srivatsaSurveyUsingLarge} conduct a comprehensive survey evaluating LLM capabilities across major declarative and procedural configuration languages, concluding that models frequently generate syntactically plausible configurations that nonetheless fail to deploy in real-world environments. This divergence stems from the models' lack of awareness of evolving best practices, deprecated configuration attributes, and fast-moving upstream dependencies.

Recent benchmark studies substantiate these findings. Kon et~al.~\cite{konIaCEvalCodeGeneration} introduce IaC-Eval, demonstrating that contemporary models like GPT-4 achieve less than a 20\% pass@1 rate when generating production cloud infrastructure programs from natural language requirements. Similarly, Zhang et~al.~\cite{zhangDeployabilityCentricInfrastructureasCodeGeneration2025} emphasize that deployability, rather than surface syntax correctness, is the decisive metric for automated system generation. A foundational vulnerability in automated environment synthesis is package hallucination~\cite{ImportingPhantomsMeasuring}. When LLMs encounter niche software requirements or less represented languages in their training corpora, they frequently fabricate non-existent package identifiers. In untrusted software pipelines, these hallucinations expose systems to package squatting security risks and break build executions, underscoring the critical need for deterministic external validation mechanisms.

\subsection{Reproducible Challenge and Lab Environments for Cybersecurity Education}

The necessity for hermetic and rapid environment deployment is particularly evident in interactive computing competitions and laboratory exercises~\cite{guptaCTFArchiveCapture2025}. In cybersecurity education, replicating Capture The Flag (CTF) challenges require students to configure their systems, resolve missing dependencies, and unclear documentation. The time spent setting up would've otherwise been better spent on understanding the underlying concept that the challenge were meant to teach. The setup overhead may also discourage learners.


To alleviate the friction of students manually and imperatively reconstructing complex exercise environments, Gupta et~al.~\cite{guptaCTFArchiveCapture2025} introduce CTF Archive, a centralized platform that preserves historical competition challenges within pre-configured, self-contained Docker container images. By eliminating the necessity for manual system configuration and dependency resolution, CTF Archive enables students to instantiate ready-to-use exercise sandboxes with minimal preparatory overhead, allowing them to focus their cognitive effort directly on conceptual understanding and vulnerability discovery. However, the authors explicitly acknowledge a fundamental pedagogical trade-off inherent in automated rehosting: by completely shielding learners from the operational deployment pipeline, students miss valuable experiential learning opportunities related to server administration, container authoring (such as composing Dockerfiles), and developing deployment automation scripts.

While container registries successfully preserve static binary snapshots, they exhibit notable pedagogical and operational drawbacks~\cite{guptaCTFArchiveCapture2025}. Container images are opaque binary blobs that obscure the granular build instructions and dependency trees from students, hindering their understanding of how software stacks are structured. Furthermore, maintaining collections of multi-gigabyte container images across dozens of academic lab modules imposes steep storage demands on university infrastructure and student personal devices alike.

\subsection{Declarative and Functional Package Management in Education}

To circumvent the heavy computational overhead of virtualized container engines, educators and researchers have begun exploring purely functional package management via Nix Flakes~\cite{5BuildingReproducible,duNixRescueReproducible2026a}. Hasanovi{\'c} et~al.~\cite{5BuildingReproducible} report on the successful migration of university C++ lab environments from bulky Docker containers to lightweight, declarative Nix Flake environments hosted on public Git repositories. Their findings highlight that Nix-based development environments provide students with rapid shell instantiation, native filesystem performance, and fine-grained dependency pinning without requiring background virtual machines or elevated root privileges.

The theoretical strengths of functional software deployment have been documented extensively in research literature~\cite{dolstraPurelyFunctionalSystema,dolstraNixOSPurelyFunctional2010}. By treating package builds as pure mathematical functions whose outputs depend exclusively on cryptographically hashed inputs, Nix guarantees that identical declarations yield identical binaries. Bandt~\cite{bandtAssessmentReproducibilityNix2025} conducts an empirical assessment of reproducibility in Nixpkgs, confirming that functional package management drastically minimizes non-deterministic build behaviors compared to traditional package registries. In scientific and academic computing, Hausch et~al.~\cite{hauschImprovingReproducibilityScientific2025} and Malka et~al.~\cite{malkaReproducibilityBuildEnvironments2024,malkaDoesFunctionalPackage2025} demonstrate that functional package declarations consistently outperform imperative shell scripts in preserving computational environments across multi-year academic timelines.

\subsection{Agentic Skills and Tool-Augmented Assistance}

While Nix Flakes offer remarkable reproducibility and lightweight execution, their declarative DSL poses an exceptionally steep learning curve for students and faculty unfamiliar with lazy functional programming~\cite{ModularityLazyEvaluation}. Recent research in agent harness architectures demonstrated that autonomous LLMs can execute system administration tasks when supplied with structured external context and tool access~\cite{mengAgentHarnessLarge2026}. Ling et~al.~\cite{lingAgentSkillsDataDriven2026} evaluate agent skills, showing that modular packages containing domain-specific templates, workflow rules, and reference metadata significantly improve an agent's task accuracy and behavioral consistency on complex, specialized workflows.

Integrating standardized tool protocols further amplifies agent reliability~\cite{borkLanguageServerProtocol2023}. By coupling LLM reasoning with real-time package search servers via the Model Context Protocol (MCP), agents gain the capacity to query current package registries dynamically, resolving the package hallucination challenges identified by earlier studies~\cite{ImportingPhantomsMeasuring}.

\section{Summary of Research Gaps}

Existing literature demonstrates a clear trade-off in educational environment management. Centralized servers introduce administrative friction and network vulnerabilities~\cite{valstarUsingDevContainersStandardize2020}. OCI DevContainers deliver isolation at the expense of excessive virtualized hardware resources and storage consumption~\cite{NixBetterDocker2024}. Nix Flakes provide the ideal balance of native performance, fine-grained reproducibility, and lightweight host execution, yet their syntax complexity restricts broad classroom adoption~\cite{5BuildingReproducible}.

Although modern LLMs excel at conversational code synthesis, they exhibit high error and hallucination rates on niche domain specific languages like Nix~\cite{ImportingPhantomsMeasuring,konIaCEvalCodeGeneration}. No prior work has investigated combining specialized custom agent skills, real-time MCP registry retrieval, and deterministic Nix evaluation feedback to bridge the cognitive divide between natural language course requirements and functional Nix Flake specifications. This thesis directly investigates this unexplored intersection.
